\section{Problem Definition}

Let $\mathcal{P}=\{p_1,\ldots,p_N\}$ denote the catalogue of whisky products extracted from the source PDF. Each product $p_i$ contains two forms of evidence: a structured text record $t_i$ and a product image $v_i$. The structured text record includes fields such as product name, brand or distillery, style, region, country, age, ABV, product tasting notes, and brand context. The product image is the cropped catalogue bottle image associated with the product. From these two sources, we construct three product representations: a text representation $p_i^{T}=t_i$, an image representation $p_i^{I}=v_i$, and a multimodal representation $p_i^{M}=(t_i,v_i)$.

We study retrieval over this fixed catalogue. A user query $q$ can be one of three modality types:
\begin{enumerate}
    \item text-only, where $q=t_q$ is a natural language product request;
    \item image-only, where $q=v_q$ is either a cropped bottle image or a whole scene image;
    \item image-text, where $q=(t_q,v_q)$ combines a whole scene image with a short textual intent, such as ``What is the whisky in this image?''.
\end{enumerate}

The task is to return a ranked list of catalogue products for each query under a selected product representation. Given an embedding function $f(\cdot)$, we represent queries and products in a shared vector space and rank products by cosine similarity:
\[
    s(q,p_i^{r})=
    \frac{f(q)^\top f(p_i^{r})}
    {\lVert f(q)\rVert_2 \lVert f(p_i^{r})\rVert_2},
\]
where $r \in \{T,I,M\}$ denotes the product representation used for retrieval. The retrieved list is obtained by sorting all products in decreasing order of $s(q,p_i^{r})$.

We instantiate this task with three modality-matched retrieval settings. In the text-only setting, text queries are compared with product text embeddings, i.e., $q=t_q$ is ranked against $p_i^{T}$. In the image-only setting, image queries are compared with product image embeddings, i.e., $q=v_q$ is ranked against $p_i^{I}$. This setting includes an ablation between cropped bottle query images and whole scene query images. In the image-text setting, whole scene images plus a textual question are compared with product embeddings built from both structured text and cropped product images, i.e., $q=(t_q,v_q)$ is ranked against $p_i^{M}$.

We additionally consider two cross-representation ablations. First, text queries are ranked against multimodal product embeddings, $q=t_q$ against $p_i^{M}$, to test whether product-side visual evidence improves or distorts text search. Second, image queries are ranked against multimodal product embeddings, $q=v_q$ against $p_i^{M}$, to test whether product-side textual evidence improves or distorts visual search. These ablations rely on the assumption that the embedding model maps text, image, and multimodal inputs into a compatible vector space; the experiment measures whether that assumption is useful in this catalogue retrieval setting.

For evaluation, each query is associated with a set of relevant catalogue product identifiers. These labels are not user preference labels; they are product-level relevance labels derived from catalogue metadata, tasting descriptions, and manually reviewed movie or television scene references. The objective is therefore to measure similarity retrieval quality, not personalized recommendation quality.
